\section{胜利之后，边疆仍要供给}
\label{sec:tang-frontiers-silk-road}

630年唐军破东突厥，次年安置降附部众；634年击吐谷浑，638年松州之战后又转入与吐蕃的婚盟。这里没有一条从胜利到服从的直线：俘获可改变草原政治，迁置则立刻提出牧地、粮食、任官与身份分类的问题；军事对峙之后的婚盟，也不是一方永久支配另一方的证明。汉文正史多以唐廷为观察中心，突厥与吐蕃的行动逻辑须从碑铭、文书和外部材料中补问。%
\footnote{东突厥、安置、吐谷浑、松州与文成公主入吐蕃见 ST-630-EASTERN-TURKS、ST-631-TURK-SETTLEMENT、ST-634-TUYUHUN、ST-638-SONGZHOU-TIBET、ST-641-WENCHENG，依据《旧唐书》《新唐书》突厥、吐谷浑、吐蕃传与《资治通鉴》。人口、贡赐和婚盟的实际条件不宜照录为精确数字。}

高昌在640年覆亡并置西州，648年攻龟兹，657年又破西突厥；这些推进把都护、羁縻与地方州县并置在天山南北。直接置州、任都护或册立首领，分别意味着不同的行政深度，不能合称为一个无缝的“西域统治”。唐军645年久攻安市未克而撤出辽东，也提醒人们远征的极限：寒冷、粮道和城防足以抵消一时的兵力优势。%
\footnote{高昌、辽东、龟兹与西突厥事件见 ST-640-GAOCHANG、ST-644-GOGURYEO-EXPEDITION、ST-645-LIAODONG-WITHDRAWAL、ST-648-KUCHA、ST-657-WESTERN-TURKS，依据《旧唐书》相关列传、《新唐书》地理及外族传和《资治通鉴》。都护辖境、兵数与战果具有行政和军功修辞。}

在东北，唐与新罗的联合先后灭百济、在白江口击败倭援军、再攻破平壤。联军的年序可以确证，却不能把新罗、倭和百济复国军写成唐朝叙事中的附属角色；《三国史记》《日本书纪》同样有本国的编纂立场。668年以后安东都护府的设置也不是高句丽社会被立刻同质接收的凭证，而是战争胜利转化为行政安排时产生的新摩擦。%
\footnote{百济、白江口与高句丽覆亡见 ST-660-BAEKJE、ST-663-BAEKGANG、ST-668-GOGURYEO-CONQUEST，依据两《唐书》东夷传、《三国史记》《日本书纪》及《资治通鉴》。战场人数、航路和地方反应在不同传统中有异。}

670年大非川战败、676年前后河湟后撤，679年东突厥复兴；武周在692年恢复安西四镇、702年置北庭，而营州契丹起事又在695--697年暴露东北的压力。边疆不是一面不断向外推的墙，而是若干需分别供给和谈判的走廊。747年、751年对南诏的两次失利，以及怛罗斯之败，同样不可被压缩为帝国“由盛转衰”的单一信号：云南山地、河中交通、盟友关系和远征补给的限制各不相同，却共同限制了中央把军令变成稳定控制的能力。%
\footnote{吐蕃、后突厥、安西北庭、契丹、南诏与怛罗斯诸事见 ST-670-DAFEICHUAN、ST-676-HEHUANG-LOSS、ST-679-ASHIDE-REBELLION、ST-692-ANXI-RECOVERY、ST-695-KHITAN-REBELLION、ST-697-KHITAN-SUPPRESSION、ST-702-BEITING、ST-747-NANZHAO、ST-751-TALAS、ST-751-NANZHAO-DEFEAT。核心为两《唐书》诸传与《资治通鉴》；吐蕃、突厥、回鹘碑铭及中亚材料是必要校正，不能以唐方失败或胜利叙事代替他方目的。}
